% RESEARCH_DRAFT_SHA256: 517714b5d19480f28a7175c8cb7623bf690654128c077f4301d06a40e08ad4a6
% Layout and source-reference wrappers expanded for the web renderer.
\documentclass[11pt]{article}
\usepackage[margin=25mm]{geometry}
\usepackage{fontspec,unicode-math}
\setmainfont{Latin Modern Roman}
\setmathfont{Latin Modern Math}
\usepackage{amsmath,amssymb,mathtools,xurl,hyperref}
\hypersetup{colorlinks=true,urlcolor=blue,linkcolor=blue}
\setlength{\parindent}{0pt}
\setlength{\parskip}{5pt}
\setlength{\emergencystretch}{3em}
\newcommand{\sref}[2]{\hyperlink{source:#1:#2}{[S#2]}}
\newcommand{\cataloguescope}[1]{\paragraph{Source question.}#1}
\newcommand{\floor}[1]{\left\lfloor#1\right\rfloor}
\begin{document}

\section*{Projective incidence complexity: the known area bound and a scale-uniformity gap}
\noindent catalogue entry; catalogue entry. Research attempt dated 11 October 2026.
\paragraph{Result in scope.}This manuscript contains a source-specific attempt and complete proofs of the stated auxiliary calculations. It makes no new solution or novelty claim for the entire target.
\section{Original problem and preserved definition}
The following definition and source list are reproduced from the repository. Any earlier inspection dates and reported literature claims in this preserved text describe that source record, not additional verification by this attempt.
% BEGIN PRESERVED DEFINITION
\subsection{Definitions and mathematical statement}

\paragraph{Shared notation.}Standard number systems are $\mathbb N,\mathbb Z,\mathbb Q,\mathbb R,\mathbb C$. Specialized notation and index ranges follow the cited source and the definitions below.


A projective plane has points and lines with incidence relations; a geometric inequality compares numerical invariants of a configuration. Effective complexity needs a specified quantitative description or representation measure.
\paragraph{Record scope and qualifications.}
The catalogue title is retained as metadata. Named certificates, contracts, bounds or frameworks in that title are not asserted as results of the original source extract.
The original statement, projective invariant, and sphere-area inequality are missing from the extraction.

Source-record limitation: The finite effective-complexity statement and sphere-area inequality referred to by the question are absent.
\paragraph{Statement recorded in the supplied source.}
Can one formulate and prove the geometric analogue of the finite effective-complexity statement for points and lines on the projective plane, together with the corresponding sphere-area inequality? \href{https://doi.org/10.4171/owr/2021/21}{[S2]}
\paragraph{Source status.}
The supplied catalogue labels this source record open; that metadata was checked against the supplied snapshot on 2026-10-11. The label is preserved for provenance and is not an independent verification of present open status or mathematical correctness. The catalogue literature assessment is dated 2026-08-22.
\subsection{Short English statement}
Develop geometric counterparts of the cited effective-complexity statement.
\subsection{Sources}
\begin{itemize}
\item[S1] Ulam AI and the UnsolvedMath Contributors, UnsolvedMath: A Curated Collection of Open Mathematics Problems, catalogue entry. Catalogue curation and metadata: CC BY 4.0. The source extract is retained; definitions, English summaries and qualifications are editorial additions.
\url{https://huggingface.co/datasets/ulamai/UnsolvedMath}.
\item[S2] Computability Theory (2022), Oberwolfach Reports, DOI 10.4171/owr/2021/21, pp. 1149--1190
\url{https://doi.org/10.4171/owr/2021/21}.
\item[S3] Oberwolfach Report 21/2021
\url{https://ems.press/content/serial-article-files/46899}.
\item Source-specific Oberwolfach Reports licences, including CC BY-SA 4.0 where applicable, are retained; see THIRD\_PARTY\_NOTICES.md.
\end{itemize}
% END PRESERVED DEFINITION
\subsection{Verbatim supplied catalogue question}
Can one formulate and prove the geometric analogue of the finite effective-complexity statement for points and lines on the projective plane, together with the corresponding sphere-area inequality?

\section{What the source already proves}
Shen's primary report, printed pages 1177--1178, states a finite-field
complexity result, then an approximation-scale lower bound
$n+s/2$, and already proves the spherical area inequality below.
Its final open question follows a discussion of why a stronger
effective-dimension statement uses a liminf across scales. Thus
reproving the area inequality is not a solution of that final
question. The catalogue does not specify a new geometric statement
with all its quantifiers. We attempt a uniform-neighbourhood
interpretation, establish its consequence, and exhibit why its
exponent and quantifier order do not yield the stronger statement.

\section{Complete rederivation of the source area inequality}
Let $\mu$ be normalized area on the unit sphere $S^2$ in
$\mathbb R^3$. Suppose measurable $A,B\subseteq S^2$ satisfy
\[
 b\in B,\quad a\cdot b=0\quad\Longrightarrow\quad a\in A.
\]
Then $\mu(B)\le\mu(A)^2$. Take independent uniform points
$\xi,\eta$ on $S^2$ and put
$\zeta=(\xi\times\eta)/\|\xi\times\eta\|$ whenever defined.
Conditional on $\xi$, both singletons $\xi,-\xi$ have zero area,
so the exceptional event $\eta=\pm\xi$ has probability zero.
The law of $(\xi,\eta)$ is invariant under every rotation, and
the cross product commutes with rotations. Thus $\zeta$ has
rotation-invariant law, which is normalized spherical area.
One elementary way to verify the last uniqueness fact is to
average any continuous function over a uniform rotation: its
value averaged at any fixed point is its spherical average,
independent of that point. Integrating this identity against any
rotation-invariant probability law yields the same average for
all continuous functions, hence the same probability measure.
Uniform rotations here can be constructed by first selecting a
uniform axis image and then a uniform angle around that image.
Now $\zeta\in B$ implies both $\xi,\eta\in A$, so independence
gives the claimed inequality. This is the source's already-known
argument, including the antipodal exceptional case; its printed
letter $N$ in the probability line is interpreted as $B$.

\section{A rigorous neighbourhood consequence}
For a nonempty subset $E$ of the sphere write
$E_\delta=\{x:\operatorname{dist}_{\rm chord}(x,E)<\delta\}$,
where $0<\delta<1$. Under the same incidence hypothesis,
\[
 \mu(B_\delta)\le\mu(A_\delta)^2.
\]
Indeed, if $b'\in B_\delta$, choose $b\in B$ with
$\|b-b'\|<\delta$. The rotation $R$ through the smaller angle
from $b$ to $b'$ in their span fixes the perpendicular axis and
satisfies $\|R-I\|_{\rm op}=\|b-b'\|<\delta$. If $a'\cdot b'=0$,
then $a=R^{-1}a'$ is perpendicular to $b$, hence in $A$, and
$\|a'-a\|<\delta$. Thus the whole great circle normal to $b'$
lies in $A_\delta$, and the preceding area inequality applies.
The open neighbourhoods are measurable even if the original sets
need only be measurable. If $B$ is empty the inequality is trivial;
otherwise $A$ contains a circle and all logarithms below are finite.

Define the volume exponents explicitly by
\[
 d_\delta(E)=2-\frac{\log\mu(E_\delta)}{\log\delta},
 \qquad \underline d(E)=\liminf_{\delta\downarrow0}d_\delta(E),
 \qquad \overline d(E)=\limsup_{\delta\downarrow0}d_\delta(E).
\]
These are the usual tubular-volume versions of lower and upper
Minkowski dimension on a two-dimensional compact surface; only
the displayed definitions are needed here. Since $\log\delta<0$,
the preceding inequality yields, with its direction checked,
\[
 d_\delta(A)\ge1+\tfrac12d_\delta(B),\quad
 \underline d(A)\ge1+\tfrac12\underline d(B),\quad
 \overline d(A)\ge1+\tfrac12\overline d(B).
\]
For the limsup assertion take a sequence attaining the right-hand
limsup; the pointwise inequality then bounds the left-hand limsup.
For liminf, use the eventual lower bound from its definition.
This does not yield the stronger expression $\min(1+s,2)$ for
$0<s<2$, nor does a volume exponent identify individual effective
Hausdorff dimensions.

\section{Why simply strengthening the area exponent fails}
Let $0<r<\pi/2$. Take $B_r$ to be the two closed spherical caps
of angular radius $r$ about the north and south poles, and let
$A_r$ be the equatorial closed band $|a_3|\le\sin r$.
Every great circle perpendicular to a vector in $B_r$ lies in
$A_r$: if $b$ makes angle at most $r$ with a pole, the length of
the projection of that pole onto $b^\perp$ is at most $\sin r$,
which bounds the absolute third coordinate of every unit vector
in $b^\perp$. Conversely every point in that band belongs to
one of these circles, by choosing the normalized projection of
the north pole onto its perpendicular plane (with an appropriate
sign and the equatorial boundary handled directly).
Spherical coordinates, for which normalized $a_3$ is uniform on
$[-1,1]$, give
\[
 \mu(A_r)=\sin r,\qquad \mu(B_r)=1-\cos r.
\]
Hence for any fixed $p>2$ and constant $C<\infty$,
$\mu(B_r)\le C\mu(A_r)^p$ fails for small $r$, since
\[
 \frac{1-\cos r}{(\sin r)^p}\sim\tfrac12 r^{2-p}\longrightarrow\infty.
\]
The asymptotic follows from $1-\cos r=2\sin^2(r/2)$ and
$\sin r/r\to1$. Thus an improved universal area power cannot
supply the missing dimension bound. For each fixed positive $r$
both sets have interior, so this family is not itself a
counterexample to a Hausdorff-dimension assertion.

\section{The unproved exchange of point and scale}
A statement producing a point at each resolution need not produce
one point good at infinitely many appropriate resolutions, much
less at all resolutions needed for a liminf. Abstractly, let
$I_m$ be pairwise disjoint nonempty subintervals of $[0,1]$ and
$a_m(x)=\boldsymbol{1}_{I_m}(x)$. Then
\[
 \sup_x\liminf_m a_m(x)=0,
 \qquad \liminf_m\sup_x a_m(x)=1.
\]
Every point belongs to at most one interval, while every interval
contains a point, proving both assertions. This is a quantifier
counterexample, not an example of actual prefix complexities.
The area proof and its neighbourhood version give no additional
coherence that would rule out the relevant escape of witnesses
across scales. Establishing the right geometric formulation and
that coherence is the remaining task.

\section{Adversarial audit and scope}
The audit checked the source attribution of the known inequality,
the null antipodal case, the sign reversal on dividing by
$\log\delta$, the difference between fixed-scale caps and a
fractal set, and the distinction between a scale-wise maximizing
point and one point controlling a liminf. No experiment or theorem
prover was used. The proved consequences are auxiliary known
mechanisms; they do not formulate and prove the source's intended
stronger geometric counterpart. The record remains unresolved.

\section{Process and contributions}
Mathematics and model-assisted adversarial self-audit: OpenAI gpt-6-astra, ultra reasoning, as dispatched for this research batch; exact serving revision not exposed. Codex (GPT-6; exact serving revision not exposed) materially assisted with orchestration and format review. No human mathematical authorship, mathematical edits, or independent human proof checking is asserted.
The coordinating Codex (GPT-6; exact serving revision not exposed) also provided a separate model-assisted written-proof review of this manuscript. This is not independent human verification.
Contribution prompt A was applied to the full definition and supplied source row, then prompt B was applied adversarially to the derived results. The item-specific audit above records the scope checks and rejected extrapolations. Local assembly and structural checks do not establish mathematical correctness. The source catalogue status is preserved in TOP\_PROBLEM and is separate from this attempt status.
\paragraph{Integration review.} Codex (GPT-6; exact serving revision not exposed) reviewed the written arguments and expanded layout and source-reference wrappers. Repository metadata and the production web renderer were checked using Python 3.13 and JavaScript with MathJax 3.2.2. This model-assisted review is not independent human proof verification; no PDF compilation or formal proof checking is asserted.

\section{Completion Estimate}
\noindent\textbf{COMPLETION: 0\%}
The proved material is a rederivation of known or auxiliary facts rather than a substantiated new advance on the original target. The final source question asks for a corresponding geometric formulation after a liminf discussion. The known area inequality and its tubular-volume consequence do not supply a stronger effective-dimension statement or coherent point witnesses across scales.
\section{References}
The complete original Sources subsection above is retained, including catalogue attribution and licence notices. The additional list records primary sources actually consulted for this attempt on 11 October 2026; uncited source links reproduced above were not all rechecked.
\begin{itemize}
\item Alexander Shen, Complexity of lines and planes: finite case, OWR21/2021, pp.1177–1178; full original passage read in cached report and primary web record checked. The area proof is attributed there to Bogdanov and Kozachinskiy. \url{https://ems.press/content/serial-article-files/46899}.
\end{itemize}
\section{Final status}
\textbf{Unresolved in this attempt.} The final source question asks for a corresponding geometric formulation after a liminf discussion. The known area inequality and its tubular-volume consequence do not supply a stronger effective-dimension statement or coherent point witnesses across scales.
\end{document}
